\documentclass{beamer}
\usepackage{beamerthemesplit}
\usepackage{multirow}
\usepackage{array}
\usepackage{hyperref}
\usepackage[T1]{fontenc}
\usepackage{inconsolata}
\usepackage{xcolor,colortbl}
\usepackage[square]{natbib}
\usepackage{listings}
%%\newcommand{\newblock}{}
\DeclareGraphicsExtensions{.pdf,.png,.jpg}
\usetheme[pageofpages=of,% String used between the current page and the
                         % total page count.
          bullet=circle,% Use circles instead of squares for bullets.
          titleline=true,% Show a line below the frame title.
          alternativetitlepage=true,% Use the fancy title page.
          ]{Torino}
\definecolor{light-green}{RGB}{144,238,144}
\makeatletter
\setbeamertemplate{footline}
{
  \leavevmode%
  \hbox{%
  \begin{beamercolorbox}[wd=.333333\paperwidth,ht=2.25ex,dp=1ex,center]{author in head/foot}%
    \usebeamerfont{author in head/foot}\insertshortauthor~~\beamer@ifempty{\insertshortinstitute}{}{(\insertshortinstitute)}
  \end{beamercolorbox}%
  \begin{beamercolorbox}[wd=.333333\paperwidth,ht=2.25ex,dp=1ex,center]{title in head/foot}%
    \usebeamerfont{title in head/foot}\insertshorttitle
  \end{beamercolorbox}%
  \begin{beamercolorbox}[wd=.333333\paperwidth,ht=2.25ex,dp=1ex,right]{date in head/foot}%
    \usebeamerfont{date in head/foot}\insertshortdate{}\hspace*{2em}
%    \insertframenumber{} / \inserttotalframenumber\hspace*{2ex} % DELETED
  \end{beamercolorbox}}%
  \vskip0pt%
}
\makeatother
\begin{document}
\author{{\bf Casey Stella}}
\institute[Hortonworks]{\includegraphics[width=40px,height=17px]{logo}}
\title{{\bf Apache Pig for Data Science}}
\date{May 12, 2014} 

\frame{\titlepage} 

\frame{\frametitle{Table of Contents}\tableofcontents} 

\section{Preliminaries}

\frame{\frametitle{Introduction}
\begin{itemize}
\item I'm a Principal Architect at Hortonworks
\item I work primarily doing Data Science in the Hadoop Ecosystem
\item Prior to this, I've spent my time and had a lot of fun
  \begin{itemize}
  \item Doing data mining on medical data at Explorys using the Hadoop ecosystem
  \item Doing signal processing on seismic data at Ion Geophysical using MapReduce
  \item Being a graduate student in the Math department at Texas A\&M in algorithmic complexity theory
  \end{itemize}
\item I'm going to talk about Apache Pig's role for doing scalable data science.
\end{itemize}
}

\subsection{Apache Pig}

\frame{\frametitle{Apache Pig: What is it?}
Pig is a high level scripting language for operating on large datasets inside Hadoop\pause
\begin{itemize}
\item Transforms high level data operations into MapReduce/Tez jobs\pause
\item Optimizes such that the minimal number of MapReduce/Tez jobs need be run\pause
\item Familiar relational primitives available \pause
\item Extensible via User Defined Functions and Loaders for customized data processing and formats
\end{itemize}
}

\begin{frame}[fragile]
\frametitle{Apache Pig: An Familiar Example}

\begin{lstlisting}
SENTENCES= load '...' as (sentence:chararray);
WORDS = foreach SENTENCES 
        generate flatten(TOKENIZE(sentence)) 
        as word;
WORD_GROUPS = group WORDS by word;
WORD_COUNTS = foreach WORD_GROUPS 
              generate group as word, COUNT(WORDS);
store WORD_COUNTS into '...';
\end{lstlisting}
\end{frame}

\section{Pig in the Data Science Toolbag}

\subsection{Understanding Your Data}

\frame{\frametitle{Understanding Data}
\begin{exampleblock}{}
  {\Large ``80\% of the work in any data project is in cleaning the data.''}
  \vskip5mm
  \hspace*\fill{\small--- D.J. Patel in {\em Data Jujitsu} }
\end{exampleblock}
}

\frame{\frametitle{Understanding Data}
A core pre-requisite to analyzing data is understanding data's shape and distribution.
This requires (among other things):
\begin{itemize}
\item Computing distribution statistics on data
\item Sampling data
\end{itemize}
}
\frame{\frametitle{Understanding Data: Datafu}
An Apache Incubating project called {\bf datafu}\footnote{http://datafu.incubator.apache.org/} 
provides some of these tooling in the form of Pig UDFs:
\begin{itemize}
\item Computing quantiles of data\pause
\item Sampling
  \begin{itemize}
  \item Bernoulli sampling by probability (built into pig)\pause
  \item Simple Random Sample \pause
  \item Reservoir sampling\pause
  \item Weighted sampling without replacement\pause
  \item Random Sample with replacement
  \end{itemize}
\end{itemize}
}

\subsection{Machine Learning with Pig}

\frame{\frametitle{What is Machine Learning?}
Machine learning is the study of systems that can learn from data.
The general tasks fall into one of two categories:\pause
\begin{itemize}
\item Unsupervised Learning
  \begin{itemize}
  \item Clustering
  \item Outlier detection
  \item Market Basket Analysis
  \end{itemize}\pause
\item Supervised Learning
  \begin{itemize}
  \item Classification
  \item Regression
  \item Recommendation
  \end{itemize}
\end{itemize}
}

\frame{\frametitle{Building Machine Learning Models with Pig}
Machine Learning at scale in Hadoop generally falls into two categories:
\begin{itemize}
\item Build one large model on all (or almost all) of the data 
\item Sample the large dataset and build the model based on that sample
\end{itemize} \pause
\vfill
Pig can assist in intelligently sampling down the large data into a training set.  You can then use your favorite
ML algorithm (which can be run on the JVM) to generate a machine learning model.
}

\subsection{Applying Models with Pig}
\frame{\frametitle{Applying Models with Pig}
Pig shines at batch application of an existing ML model.  This generally is of the form:
\vfill
\begin{itemize}
\item Train a model out-of-band\vfill \pause
\item Write a UDF in Java or another JVM language which can apply the model to a data point\vfill\pause
\item Call the UDF from a pig script to distribute the application of the model across your data in parallel\vfill
\end{itemize}
}


\section{Unstructured Data Analysis with Pig}

\frame{\frametitle{What is Natural Language Processing?}
\begin{itemize}
\item Natural language processing is the field of Computer Science, Linguistics \& Math that covers computer understanding and manipulation of human language.
  \begin{itemize}
  \item Historically, linguists hand-coded rules to accomplish much analysis
  \item Most modern approaches involves using Machine Learning
  \end{itemize}\pause
\vfill
\item Mature field with many useful libraries on the JVM
  \begin{itemize}
  \item Apache OpenNLP
  \item Stanford CoreNLP
  \item MALLET
  \end{itemize}
\end{itemize}
\vfill
}

\frame{\frametitle{Natural Language Processing with Large Data}
\begin{itemize}
\item Generally low-volume, complex analysis
  \begin{itemize}
  \item Big companies often don't have a ton of natural language data
  \item Dropped previously because they were unable to analyze
  \end{itemize}\pause
\vfill
\item Sometimes high-volume, complex analysis
  \begin{itemize}
  \item Search Engines
  \item Social media content analysis
  \end{itemize}\pause
\vfill
\item Typically many small-data problems in parallel
  \begin{itemize}
  \item Often requires only the context of a single document
  \item Ideal for encapsulating as Pig UDFs
  \end{itemize}
\end{itemize}
\vfill
}

\frame{\frametitle{Natural Language Processing: Demo}

\begin{itemize}
\item Stanford CoreNLP integrated the work of Richard Socher, et al ~\citep{Socher-etal:2013} using recursive deep neural networks to
predict sentiment of movie reviews.\vfill
\item There is a large set of IMDB movie reviews  used to analyze sentiment analysis~\citep{maas-EtAl}.\vfill
\item Let's look at how to encapsulate this into a Pig UDF and run on some movie review data.
\vfill
\end{itemize}
}

\frame{\frametitle{Results}
\begin{itemize}
\item Executing on a sample of size $1022$ Positive and Negative documents.
\item Overall Accuracy of $77.2\%$
\end{itemize}

\begin{tabular}{l|l|c|c|c}
\multicolumn{2}{c}{}&\multicolumn{2}{c}{Actual}&\\
\cline{3-4}
\multicolumn{2}{c|}{}&Positive&Negative&\multicolumn{1}{c}{Total}\\
\cline{2-4}
\multirow{2}{*}{Predicted}& Positive & $367$ & $114$ & $481$\\
\cline{2-4}
& Negative & $119$ & $422$ & $541$\\
\cline{2-4}
\multicolumn{1}{c}{} & \multicolumn{1}{c}{Total} & \multicolumn{1}{c}{$486$} & \multicolumn{1}{c}{$536$} & \multicolumn{1}{c}{$1022$}\\
\end{tabular}
}

\section{Questions \& Bibliography}

\frame{\frametitle{Questions}
Thanks for your attention!  Questions? 
\vfill
\begin{itemize}
\item Code \& scripts for this talk available on my github presentation page.\footnote{http://github.com/cestella/presentations/}
\item Find me at http://caseystella.com 
\item Twitter handle: @casey\_stella 
\item Email address: cstella@hortonworks.com
\end{itemize}
}

\frame{\frametitle{Bibliography}
\bibliography{Pig_for_DataScience}{}
\bibliographystyle{plain}
}
\end{document}
